\documentclass[12pt]{article}

% Layout.
\usepackage[top=1.2in, bottom=0.75in, left=1in, right=1in, headheight=1.0in, headsep=0pt]{geometry}

% Fonts.
\usepackage{mathptmx}
\usepackage[scaled=0.86]{helvet}
\renewcommand{\emph}[1]{\textsf{\textbf{#1}}}

% TiKZ.
\usepackage{tikz, pgfplots}
\usetikzlibrary{calc}
%\pgfplotsset{my style/.append style={axis x line=middle, axis y line=middle, xlabel={$x$}, ylabel={$y$}}}
%\pgfplotsset{compat=1.16}

% Misc packages.
\usepackage{amsmath,amssymb,latexsym}
\usepackage{graphicx}
\usepackage{array}
\usepackage{xcolor}
\usepackage{multicol}

% Commands to set various header/footer components.
\makeatletter
\def\doctitle#1{\gdef\@doctitle{#1}}
\doctitle{Use {\tt\textbackslash doctitle\{MY LABEL\}}.}
\def\docdate#1{\gdef\@docdate{#1}}
\docdate{Use {\tt\textbackslash docdate\{MY DATE\}}.}
\def\doccourse#1{\gdef\@doccourse{#1}}
\let\@doccourse\@empty
\def\docscoring#1{\gdef\@docscoring{#1}}
\let\@docscoring\@empty
\def\docversion#1{\gdef\@docversion{#1}}
\let\@docversion\@empty
\makeatother

% Headers and footers layout.
\makeatletter
\usepackage{fancyhdr}
\pagestyle{fancy}
\fancyhf{} % Clears all headers/footers.
\lhead{\emph{\@doctitle\hfill\@docdate}
\ifnum \value{page} > 1\relax\else\\
\emph{Name: \rule{3.5in}{1pt}\hfill \@docscoring}
\fi}

\rfoot{\emph{\@docversion}}
\lfoot{\emph{\@doccourse}}
\cfoot{\emph{\thepage}}
\renewcommand{\headrulewidth}{0pt}%
\makeatother

% Paragraph spacing
\parindent 0pt
\parskip 6pt plus 1pt

% A problem is a section-like command. Use \problem{5} for a problem worth 5 points.
\newcounter{probcount}
\newcounter{subprobcount}
\setcounter{probcount}{0}
\newcommand{\problem}[1]{%
\par
\addvspace{4pt}%
\setcounter{subprobcount}{0}%
\stepcounter{probcount}%
\makebox[0pt][r]{\emph{\arabic{probcount}.}\hskip1ex}\emph{[#1 points]}\hskip1ex}
\newcommand{\thesubproblem}{\emph{\alph{subprobcount}.}}

% Subproblems are an enumerate-like environment with a consistent
% numbering scheme. 
% Use \begin{subproblems}\item...\item...\end{subproblems}
\newenvironment{subproblems}{%
\begin{enumerate}%
\setcounter{enumi}{\value{subprobcount}}%
\renewcommand{\theenumi}{\emph{\alph{enumi}}}}%
{\setcounter{subprobcount}{\value{enumi}}\end{enumerate}}

% Blanks for answers in normal and math mode.
\newcommand{\blank}[1]{\rule{#1}{0.75pt}}
\newcommand{\mblank}[1]{\underline{\hspace{#1}}}
\def\emptybox(#1,#2){\framebox{\parbox[c][#2]{#1}{\rule{0pt}{0pt}}}}

% Misc.
\renewcommand{\d}{\displaystyle}
\newcommand{\ds}{\displaystyle}


\doctitle{Math 252: Quiz 1}
\docdate{Fall 2025}
\doccourse{}
\docversion{}
\docscoring{\fbox{{\LARGE \strut}\blank{0.8in} / 24}}

\begin{document}
24 points possible; each part is worth 2 points.  No aids (book, notes, calculator, phone, etc.) are permitted.  Show all work and use proper notation for full credit.  Answers should be in reasonably-simplified form.

\problem{12}  Compute the derivatives of the following functions.


\begin{subproblems}
 % Like quiz1 c
\item   $\ds f(x) =\cos(ax^2+b)$ for fixed constants $b$ and $c$
\vfill

%basic rules they were told in advance they needed to know
\item  $\ds f(\theta)=\arcsin(\theta) -2^\theta + \frac{1}{e}$
\vfill



%easy product %natural log
\item  $\ds H(t) = t^{2.4} \ln(2t+1)$

\vfill

\newpage
\quad \\
% two simple chain rules
\item   $\ds f(x) =5e^{x/2}+\sin^2(x)$
\vfill

% quotient rule
\item  $\ds h(x) = \frac{x+\sin(x)}{x+\pi}$
\vfill

%
\item  $\ds h(x) = \frac{1}{5x}+ \frac{\sqrt{x}}{10}$

\vfill
\end{subproblems}


\newpage
\problem{12}  Compute the following antiderivatives (indefinite integrals) and definite integrals.  Remember that antiderivatives need a ``$+C$''.


\begin{subproblems}

%easy start
\item $\ds \int \sin(2x) + \sqrt{x}\; dx$
\vfill


%% HW 1.5 #274
\item $\ds \int (6x-5)^{(1/4)}  \; dx$
\vfill


%% like practice quiz 1 #f
\item $\ds \int \frac{\sec^2(x)}{\tan(x)} \; dx$
\vfill

\newpage
\quad


%simplify first
\item $\ds \int \frac{2x^3+x^2 -12}{x^2}\; dx$
\vfill

%% HW 1.5 # 294
\item $\ds \int_{-1}^0 \frac{x}{(2+x)^2} \; dx$

\vfill

%% HW 1.7 # 424 
\item $\ds \int \frac{e^t}{1+e^{(2t)}} \; dt$
\vfill

\end{subproblems}
\end{document}